%%=============================================================================
%% Methodologie
%%=============================================================================

\chapter{\IfLanguageName{dutch}{Methodologie}{Methodology}}%
\label{ch:methodologie}

Dit hoofdstuk beschrijft de methodologie die gebruikt wordt om de onderzoeksvragen te beantwoorden. Het onderzoek combineert een literatuurstudie (Hoofdstuk~\ref{ch:stand-van-zaken}) met een praktijkgerichte casestudy bij Citymesh waarin een implementatie van NetBox Discovery wordt opgezet en geëvalueerd op concrete sites.

\section{Onderzoeksopzet}

Dit onderzoek volgt een \emph{praktijkgerichte casestudy-opzet} met een combinatie van:
\begin{itemize}
  \item \emph{kwalitatieve dataverzameling}: gesprekken met netwerkbeheerders, observaties van de bestaande werkwijze en analyse van implementatielogs (zie \S\ref{sec:kwalitatieve-data});
  \item \emph{gerichte voor/na-metingen}: correctheid van device-attributen en kabelrelaties in NetBox vóór en na de toepassing van NetBox Discovery, op vier sites met verschillende profielen (zie \S\ref{sec:operationalisering}).
\end{itemize}

De keuze voor een casestudy is verantwoord omdat de onderzoeksvraag contextafhankelijk is: de bruikbaarheid van automatische discovery hangt af van de bestaande processen, de toolketen en de vendor-mix binnen Citymesh. Het doel is niet om een universele effectgrootte te claimen, maar om in een realistische operationele setting te bepalen welke verbetering haalbaar is, waar de grenzen liggen en welke procesaanpassingen nodig zijn voor duurzame borging.

\section{Kwalitatieve dataverzameling}
\label{sec:kwalitatieve-data}

Om de huidige werkwijze, pijnpunten en verwachtingen rond netwerkdocumentatie bij Citymesh in beeld te brengen, werden in de loop van het onderzoek meerdere gesprekken gevoerd met de \emph{Network Engineer} en \emph{cosupervisor} bij Citymesh, Tom Pirlet, en met andere netwerkbeheerders die in de praktijk met NetBox werken. De gesprekken waren \emph{semi-gestructureerd}: er werd vertrokken van een vooraf bepaalde lijst van thema's en kernvragen, maar er was bewust ruimte om door te vragen op concrete voorbeelden uit de operationele praktijk. Naast deze gesprekken werden ook observaties gedaan tijdens het uitvoeren van imports en het inspecteren van bestaande NetBox-data, en werd de output van implementatielogs (Orb agent, Diode, NetBox API) gebruikt om de bevindingen te onderbouwen.

\subsection{Thema's en gespreksonderwerpen}

De gesprekken werden opgebouwd rond vijf thema's, met telkens enkele leidende vragen.

\paragraph{Thema 1 --- Huidige documentatieworkflow.}
\begin{itemize}
  \item Hoe wordt een nieuw site, switch of access point op dit moment in NetBox gedocumenteerd?
  \item Welke tools (Excel, NetBox-import, scripts, monitoring) komen daarbij in welke volgorde tussen?
  \item Wie is verantwoordelijk voor het bijwerken van NetBox bij wijzigingen (vervangingen, uitbreidingen, decommissions)?
\end{itemize}

\paragraph{Thema 2 --- Pijnpunten en frequentie ervan.}
\begin{itemize}
  \item Welke types fouten in NetBox komen in de praktijk het vaakst voor (verkeerde naam, verkeerd type, ontbrekende kabels, foute IP)?
  \item Bij welke gebeurtenissen treden die fouten typisch op (initiële opbouw, vervanging, verhuizing, tijdelijke setup)?
  \item Hoeveel tijd gaat er bij benadering naar het corrigeren van een typische site, en wat is de impact wanneer NetBox achterloopt?
\end{itemize}

\paragraph{Thema 3 --- Naamgevings- en data\-conventies.}
\begin{itemize}
  \item Welke naamgevingsconventies worden gebruikt voor sites, switches en access points?
  \item Hoe worden IP-adressen toegekend (vaste ranges per site, opeenvolgend per AP-reeks)?
  \item Welke velden in NetBox zijn ``contractueel'' verplicht en welke worden in de praktijk vaak overgeslagen?
\end{itemize}

\paragraph{Thema 4 --- Verwachtingen ten aanzien van automatische discovery.}
\begin{itemize}
  \item Welke datatypes zijn belangrijk om automatisch te kunnen ontdekken (devices, interfaces, IP, kabels, LLDP-buren)?
  \item Welke vendors moeten minstens ondersteund worden (MikroTik, Huawei, RUCKUS)?
  \item Hoe diep moet de tool ingrijpen op bestaande data (alleen toevoegen, of ook bijwerken/verwijderen bij drift)?
  \item Welke randvoorwaarden zijn er rond credentials, secrets en netwerktoegang per site?
\end{itemize}

\paragraph{Thema 5 --- Organisatorische inbedding.}
\begin{itemize}
  \item Hoe vaak zou een discoveryflow per site moeten draaien om voldoende actueel te blijven?
  \item Wie zou alarmen of drift-meldingen moeten oppikken?
  \item Welke aanpassingen aan procedures (cleanup vóór discovery, validatieronde na discovery) zijn realistisch?
\end{itemize}

\subsection{Wat de gesprekken concreet hebben opgeleverd}

De gesprekken hebben drie soorten input opgeleverd voor dit onderzoek. Ten eerste werd het beeld van de \emph{pijnpunten} gevalideerd en aangescherpt; de vaststelling dat een NetBox-record na een AP-vervanging vaak niet wordt bijgewerkt, kwam meermaals expliciet ter sprake en wordt verder uitgewerkt in Hoofdstuk~\ref{ch:context}. Ten tweede werden de \emph{requirements} voor de oplossing concreter, onder meer rond vendor-dekking (MikroTik, Huawei, RUCKUS) en rond de noodzaak van een \emph{cleanup}-stap vóór herdiscovery. Ten derde werden enkele \emph{operationele randvoorwaarden} duidelijk, zoals het gebruik van OAuth2-credentials, de noodzaak om geen plain-text wachtwoorden in YAML-bestanden te zetten en het feit dat agents per site moeten kunnen draaien. Deze input is verwerkt in de implementatie (Hoofdstuk~\ref{ch:implementatie}) en in de aanbevelingen voor Citymesh (Hoofdstuk~\ref{ch:analyse}).

\section{Operationalisering van de kernindicatoren}
\label{sec:operationalisering}

Om de centrale onderzoeksvraag toetsbaar te maken, werden de volgende kernindicatoren vooraf vastgelegd:
\begin{itemize}
  \item \emph{Actualiteit}: mate waarin NetBox-data overeenkomt met de actuele infrastructuurtoestand (geverifieerd via discovery-output);
  \item \emph{Correctheid van inventaris}: juistheid van device-attributen, in het bijzonder \emph{naam} en \emph{device type}, gemeten als ``aantal correct / totaal'';
  \item \emph{Volledigheid van topologie}: aanwezigheid en correctheid van kabelrelaties tussen switches onderling, tussen switches en routers en tussen access points en hun upstream-switch, opnieuw gemeten als ``aantal correct / totaal'';
  \item \emph{Resterende manuele last}: datavelden of relaties die ook na discovery nog manueel of via aanvullende verwerking met \texttt{NetworkTool} verwerkt moeten worden.
\end{itemize}

Deze indicatoren worden in de casestudy (Hoofdstuk~\ref{ch:casestudy}) per site beoordeeld via een voor/na-vergelijking. De casestudy omvat vier sites met opzettelijk verschillende profielen, zodat de bevindingen niet afhankelijk zijn van \'e\'en specifieke configuratie:

\begin{itemize}
  \item \emph{Ridefort}: kleinere site (1 router, 1 switch, 6 AP's) waarvan de NetBox-documentatie uit 2020 dateert en waar in 2022 een AP-vervanging gebeurde. Primaire voor/na-casus.
  \item \emph{Deleers Go4tech}: grotere site (1 router, 1 firewall, 11 switches, 48 AP's) met relatief recente documentatie (\textasciitilde{}2 maanden oud). Gebruikt om de waarde van \emph{periodieke} discovery te toetsen op recent gedocumenteerde data.
  \item \emph{Koninklijke Stallingen}: kleinere site (1 router, 2 switches, 5 AP's) met al consistente naamgeving maar zonder gedocumenteerde kabelrelaties.
  \item \emph{WZC Riethove}: grote site (1 router, 4 switches, 70 AP's) met consistente naamgeving en zonder gedocumenteerde kabelrelaties; ook bedoeld om het effect van een fysiek offline switch op de meting te observeren.
\end{itemize}

\section{Onderzoeksfasen en planning}

Het onderzoek is opgedeeld in zes fasen die deels sequentieel en deels parallel werden uitgevoerd. Tabel~\ref{tab:fasen-planning} vat de fasen, hun deliverables en de tijdschattingen samen.

\begin{table}[H]
  \centering
  \caption[Onderzoeksfasen, deliverables en tijdschatting]{Onderzoeksfasen, deliverables en tijdschatting}
  \label{tab:fasen-planning}
  \footnotesize
  \begin{tabular}{p{0.04\linewidth} p{0.30\linewidth} p{0.50\linewidth} p{0.06\linewidth}}
    \toprule
    \textbf{Fase} & \textbf{Inhoud} & \textbf{Deliverable} & \textbf{Duur} \\
    \midrule
    1 & Literatuurstudie en kwalitatieve dataverzameling & Hoofdstuk~\ref{ch:stand-van-zaken} en gevalideerde lijst van pijnpunten/requirements & 2 wk \\
    2 & Toolvergelijking & Vergelijkende bespreking (geïntegreerd in \S\ref{sec:alternatieven}) en motivering van de keuze voor de NetBox Discovery-stack & 1 wk \\
    3 & Baselinemetingen & Per site een ``voor''-meting in NetBox (zie Hoofdstuk~\ref{ch:casestudy}) & 1 wk \\
    4 & Implementatie en configuratie & Werkende Orb-agent + Diode-keten en \texttt{NetworkTool} (Hoofdstuk~\ref{ch:implementatie}) & 2 wk \\
    5 & Evaluatie en metingen & Per site een ``na''-meting met voor/na-tabellen (Hoofdstuk~\ref{ch:casestudy}) & 3--4 wk \\
    6 & Analyse, aanbevelingen en scriptie & Hoofdstuk~\ref{ch:analyse} en Hoofdstuk~\ref{ch:conclusie} & 3--4 wk \\
    \bottomrule
  \end{tabular}
\end{table}

\section{Tools en technologieën}

Voor de uitvoering van dit onderzoek werden onder meer de volgende tools gebruikt. De achterliggende werking en alternatieven worden besproken in Hoofdstuk~\ref{ch:stand-van-zaken}.

\begin{itemize}
  \item \emph{NetBox} als bestaande SSOT bij Citymesh, gebruikt voor zowel baseline- als evaluatiemetingen\autocite{NetBoxOSS}.
  \item \emph{NetBox Discovery, Diode en de Orb agent} als primaire automatische detectiestack\autocite{NetBoxDiscoveryDocs,DiodeDocs,NetBoxLabsOrbAgent}.
  \item \texttt{NetworkTool}: een eigen ontwikkeling, opgebouwd als Streamlit-applicatie\autocite{Streamlit}, die de NetBox-API aanstuurt om datatypes te capteren die de standaard discoveryflow niet automatisch in NetBox plaatst (in het bijzonder LLDP-gebaseerde kabelrelaties en de RUCKUS vSZ-koppeling\autocite{RuckusVSZ}). NAPALM\autocite{NAPALM} wordt gebruikt als multi-vendor abstractielaag.
  \item \emph{Docker} om de Orb agent en Diode reproduceerbaar uit te rollen per site.
  \item \emph{Python en de NetBox REST/GraphQL-API} voor analyse-, sync- en validatiescripts.
\end{itemize}

\section{Kwaliteitsborging en beperkingen van de methode}

Om de betrouwbaarheid van de bevindingen te verhogen, zijn drie principes toegepast:
\begin{itemize}
  \item \emph{Triangulatie van bronnen}: combinatie van NetBox-observaties, kwalitatieve input uit gesprekken met netwerkbeheerders en implementatielogs van Orb/Diode/NetworkTool.
  \item \emph{Transparante voor/na-logica}: baseline en eindmeting worden op dezelfde indicatoren vergeleken, met dezelfde meetdefinities (zie Hoofdstuk~\ref{ch:casestudy}).
  \item \emph{Expliciete afbakening}: claims worden beperkt tot de geobserveerde scope (vier sites, geselecteerde vendors, NetBox als SSOT) en tot indicatoren die direct in NetBox toetsbaar zijn.
\end{itemize}

De belangrijkste methodologische beperkingen zijn de beperkte steekproefomvang (vier sites), de afhankelijkheid van beschikbare discovery-protocollen per vendor en het feit dat niet elk datatype in elke context even goed meetbaar is met dezelfde automatiseringsgraad. De conclusies zijn daardoor in de eerste plaats geldig voor de Citymesh-context en moeten bij overname naar andere organisaties opnieuw lokaal gevalideerd worden.
